% !PRINTED 295
\weberpage{295}
%<<CONT
und führt zu den Kongruenzen:
\weq{5}{\begin{aligned}
1.&\quad \alpha'\alpha+\beta'\gamma \equiv \gamma'\\
2.&\quad \alpha'\beta+\beta'\delta \equiv \delta'\\
3.&\quad \gamma'\alpha+\delta'\gamma \equiv -\alpha'+\gamma'\xi\\
4.&\quad \gamma'\beta+\delta'\delta \equiv -\beta'+\delta'\xi
\end{aligned}\ \ (\mathrm{mod}\ p),}
und aus 1. und 2. folgt noch:
\weq{6}{\begin{aligned}
5.&\quad \alpha' \equiv \gamma'\delta-\delta'\gamma\\
6.&\quad \beta' \equiv -\gamma'\beta+\delta'\alpha
\end{aligned}\ \ (\mathrm{mod}\ p).}
Setzt man diese Werte in (5) 3. und 4. ein, so folgt:
\[
\begin{aligned}
\gamma'(\alpha+\delta-\xi) &\equiv 0\\
\delta'(\alpha+\delta-\xi) &\equiv 0
\end{aligned}\ \ (\mathrm{mod}\ p),
\]
woraus, da $\gamma'$, $\delta'$ nicht beide durch $p$ teilbar sein können, folgt:
\weq{7}{\xi \equiv \alpha+\delta \quad (\mathrm{mod}\ p);}
ist $\xi$ so bestimmt, so folgen in (5) die Kongruenzen 3., 4. aus
1., 2. Setzt man aber $\gamma'$, $\delta'$ aus 1., 2. in
\weq{8}{\alpha'\delta'-\beta'\gamma' \equiv 1 \quad (\mathrm{mod}\ p)}
ein, so erhält man
\weq{9}{\alpha'^2\beta+\alpha'\beta'(\delta-\alpha)-\beta'^2\gamma \equiv 1 \quad (\mathrm{mod}\ p).}
Ist hieraus $\alpha'$, $\beta'$ bestimmt, so sind alle in (5), und also
auch in (2), (4) ausgesprochenen Forderungen befriedigt.

Es handelt sich also noch darum, nachzuweisen, daß die
Kongruenz (9) immer lösbar ist.

Diese Möglichkeit ist evident, wenn $\beta$ oder $-\gamma$ quadratischer
Rest von $p$ ist; denn dann genügt
\[
\begin{aligned}
\alpha' &\equiv \sqrt{\beta^{-1}}, & \beta' &\equiv 0\\
\text{oder}\quad \alpha' &\equiv 0, & \beta' &\equiv \sqrt{-\gamma^{-1}}
\end{aligned}\ \ (\mathrm{mod}\ p)
\]
der gestellten Forderung.

Im weiteren ist nun zu unterscheiden, ob die Substitution $S$
zu Klasse I, II, III des vorigen Paragraphen gehört.

Ist zunächst
\[
\text{I.} \qquad \left(\frac{\alpha+\delta}{2}\right)^2 \equiv 1 \quad (\mathrm{mod}\ p),
\]
so geht, wenn wir $S^m$ an Stelle von $S$ setzen, nach \S\,79, (9),
$\beta$ in $m\beta$, $\gamma$ in $m\gamma$ über, und es läßt sich $m$ immer so bestimmen,
daß $m\beta$ oder $-m\gamma$ quadratischer Rest von $p$ wird.

Hierher gehört auch der Fall, daß $\beta \equiv 0$ oder $\gamma \equiv 0$ und
$\alpha \equiv \delta \ (\mathrm{mod}\ p)$ ist.
